% Irodov Part III §3.4 problems 189--200, source PDF pp.128--129,
% printed pp.119--120. The user-provided book, not invented examples.
\electrocurrentproblem{189}{119}
电阻为 $R$ 的线圈通过的总电荷为 $q$，电流逐渐降为零。
求焦耳热：(a) 电流在 $\Delta t$ 内线性减小；
(b) 电流连续衰减，每经过 $\Delta t$ 减半。
\exsolution
(a) 设初始电流为 $I_0$，
$I(t)=I_0(1-t/\Delta t)$，$0\le t\le\Delta t$。
通过的电荷
$q=\int_0^{\Delta t}I\,\mathrm dt=I_0\Delta t/2$，
故 $I_0=2q/\Delta t$。
焦耳热
\[
Q_a=R\int_0^{\Delta t}I(t)^2\,\mathrm dt
=RI_0^2\frac{\Delta t}{3}
=\boxed{\frac{4Rq^2}{3\Delta t}}.
\]
(b) 令
$I(t)=I_0e^{-\lambda t}$，
每隔 $\Delta t$ 减半意味着
$\lambda=\ln2/\Delta t$。
由 $q=I_0/\lambda$，得 $I_0=q\lambda$。
于是
\[
Q_b=R\int_0^\infty I_0^2e^{-2\lambda t}\,\mathrm dt
=\frac{RI_0^2}{2\lambda}
=\boxed{\frac{Rq^2\ln2}{2\Delta t}}.
\]
\excheck 量纲均为焦耳。
第(b)问“每隔固定时段减半”表示连续指数规律，
不能擅自改为逐段阶梯电流。

\electrocurrentproblem{190}{119}
一只内阻为 $R_0$ 的直流电源，
接到图3.52中三个相同电阻 $R$ 组成的负载。
求电路产生热功率最大时 $R$ 的取值。
\exsolution
检查图3.52的导线短接节点：
第一只、第三只以及第二只 $R$
实际上都跨在同一对等势节点上，
故三只电阻并联，
$R_{\rm L}=R/3$。
若电动势为 $\mathcal E$，
电路热功率
$P_{\rm total}=I^2(R_0+R_{\rm L})
=\mathcal E^2/(R_0+R/3)$
会随 $R$ 变小单调增加，\emph{没有正的内部极大值}。
但若题目所称“产生热功率最大”是
\emph{负载上释放的热功率}，
\[
P_{\rm L}=\frac{\mathcal E^2(R/3)}
{(R_0+R/3)^2},
\]
对负载电阻微分，
最大功率传输条件
$R_{\rm L}=R_0$，因此
\[
\boxed{R=3R_0}.
\]
\excheck 本题书末答案为 $3R_0$，
故所指的“thermal power generated in this circuit”
实际采用\emph{负载功率}口径；
已与全电路总热功率明确区分。

\electrocurrentproblem{191}{119}
证明两个并联电阻 $R_1,R_2$
在给定总电流 $I$ 时，
实际稳态分流使焦耳热功率最小。
\exsolution
设分流为 $I_1$ 与 $I_2=I-I_1$。
总焦耳功率
\[
P(I_1)=I_1^2R_1+(I-I_1)^2R_2.
\]
求导并令零：
$P'(I_1)=2I_1R_1-2(I-I_1)R_2=0$，
故
\[
\boxed{I_1=\frac{R_2}{R_1+R_2}I,\quad
I_2=\frac{R_1}{R_1+R_2}I}.
\]
二阶导数
$P''=2(R_1+R_2)>0$，
为严格极小值，
且
$P_{\min}=I^2R_1R_2/(R_1+R_2)$。
\excheck 极小条件等价于
$I_1R_1=I_2R_2$，
即两并联支路电压相同。

\electrocurrentproblem{192}{119}
蓄电池电动势 $\mathcal E=2.6\,\mathrm V$，
向外电阻输出电流 $I=1.0\,\mathrm A$，
端电压 $V=2.0\,\mathrm V$。
求电池内部产生的热功率，
以及\emph{电场力在电池内部}做功功率。
\exsolution
电池内部电压损失
$\mathcal E-V=Ir$，
故内阻 $r=(2.6-2.0)/1=0.60\,\Omega$。
电池内焦耳热
\[
\boxed{P_{\rm heat}=I^2r=I(\mathcal E-V)
=0.60\,\mathrm W}.
\]
电池内部电场沿电位\emph{降低}的方向做功；
电池内部电荷被非静电化学力逆电场推进，
电场力做功为负：
\[
\boxed{P_{\rm electric,inside}=-IV=-2.0\,\mathrm W}.
\]
\excheck 电动势贡献的非电场力功率
$\mathcal EI=2.6\,\mathrm W$，
与负的内部静电场功率及正的发热功率
满足能量平衡。不能把 $IV$ 误标为内部电场力的正功。

\electrocurrentproblem{193}{119}
直流电动机端电压固定为 $V$，
电枢绕组电阻为 $R$。
求电机有用机械功率最大时的电流、
最大有用功率及效率。
\exsolution
输入电功率 $P_{\rm in}=VI$，
电枢焦耳损耗 $P_{\rm heat}=I^2R$。
忽略其他损失，有用功率
\[
P_{\rm use}=VI-I^2R.
\]
对 $I$ 求导：
$\mathrm dP_{\rm use}/\mathrm dI=V-2IR=0$，
故
\[
\boxed{I_*=\frac V{2R}}.
\]
因二阶导数 $-2R<0$，
这是最大值；
\[
\boxed{P_{\max}=\frac{V^2}{4R}},
\qquad
\boxed{\eta_*=\frac{P_{\max}}{VI_*}=\frac12}.
\]
\excheck 50\%效率只是此理想模型的
\emph{最大输出功率工况}，
不是一般直流电机的最高效率。

\electrocurrentproblem{194}{119}
灯丝蒸发使其直径变小，
为维持原温度需将电压提高
$\eta=1.0\%$。
若灯丝向环境失热功率正比于表面积，
求灯丝直径减少的百分比。
\exsolution
温度不变时电阻率近似不变，
长度保持不变，电阻 $R\propto d^{-2}$。
电加热功率
$P_{\rm el}=V^2/R\propto V^2d^2$；
热散失 $P_{\rm loss}\propto d$。
维持温度意味着
$V^2d^2\propto d$，故 $V^2d=\text{常数}$。
若 $V_f=V_i(1+\eta)$，
\[
\frac{d_f}{d_i}=\frac1{(1+\eta)^2}.
\]
因此直径相对减少量
\[
\boxed{\frac{d_i-d_f}{d_i}
=1-(1+\eta)^{-2}\simeq2\eta=2.0\%}.
\]
\excheck 一阶近似的 2\% 与
精确值 $1-1/1.01^2\simeq1.97\%$ 一致。

\electrocurrentproblem{195}{119}
电阻为 $R$、总热容为 $C$ 的均匀导体
从初温 $T_0$ 起在 $t=0$ 接恒压 $V$。
向环境散热功率为
$k(T-T_0)$。求 $T(t)$。
\exsolution
热容乘温升速率等于电加热减散热：
\[
C\frac{\mathrm dT}{\mathrm dt}
=\frac{V^2}R-k(T-T_0).
\]
令 $\theta=T-T_0$，
则
$\dot\theta+(k/C)\theta=V^2/(RC)$。
初值 $\theta(0)=0$，
由常系数微分方程得
\[
\boxed{T(t)=T_0+
\frac{V^2}{kR}\left(1-e^{-kt/C}\right)}.
\]
\excheck $t\to\infty$ 时温升
$V^2/(kR)$，正好令电加热和环境散热相等；
书后文本若丢失负指数应以微分方程为准。

\electrocurrentproblem{196}{119--120}
图3.53中电压 $V_{AB}$ 恒定，
串联 $R_1=20\,\Omega$，
之后接并联 $R_2=30\,\Omega$
及可变 $R_x$。求令 $R_x$ 产生热功率
对其小变化一阶不敏感的 $R_x$。
\exsolution
从 $R_x$ 端看其余网络的戴维南电阻
是
$R_{\rm th}=R_1\parallel R_2
=R_1R_2/(R_1+R_2)$。
开路戴维南电压为
$V_{\rm th}=V_{AB}R_2/(R_1+R_2)$。
因此
\[
P_x=\frac{V_{\rm th}^2R_x}
{(R_{\rm th}+R_x)^2}.
\]
求导有
$\mathrm dP_x/\mathrm dR_x\propto
R_{\rm th}-R_x$，
故
\[
\boxed{R_x=R_{\rm th}
=\frac{R_1R_2}{R_1+R_2}
=\frac{20\cdot30}{20+30}=12\,\Omega}.
\]
\excheck 该点实际上是 $P_x$ 的局部极大点，
因此对于小的 $\delta R_x$，
功率的一阶改变量为零。

\electrocurrentproblem{197}{120}
图3.54中两只电源电动势
$\mathcal E_1,\mathcal E_2$，内阻
$R_1,R_2$，与电阻 $R$ 并联。
求电阻 $R$ 的热功率最大时
$R$ 的数值和最大功率。
\exsolution
从负载 $R$ 端向两并联源看，
它们的戴维南电动势为
\[
\mathcal E_{\rm th}=
\frac{\mathcal E_1R_2+\mathcal E_2R_1}{R_1+R_2},
\]
等效内阻
$R_{\rm th}=R_1R_2/(R_1+R_2)$。
于是负载功率
$P_R=\mathcal E_{\rm th}^2R/(R+R_{\rm th})^2$。
最大功率传输条件为
\[
\boxed{R_*=\frac{R_1R_2}{R_1+R_2}},
\]
\[
\boxed{P_{\max}=
\frac{\mathcal E_{\rm th}^2}{4R_{\rm th}}
=\frac{(\mathcal E_1R_2+\mathcal E_2R_1)^2}
{4R_1R_2(R_1+R_2)}}.
\]
\excheck 本题求电阻 $R$ 的功率，
不是包括两个内阻的全网络耗散。

\electrocurrentproblem{198}{120}
共有 $N=300$ 节相同电池，
每节内阻 $r=0.3\,\Omega$，
外接负载 $R=10\,\Omega$。
把电池分为 $n$ 组并联，
每组有 $N/n$ 节串联。
求使负载热功率最大的并联组数 $n$。
\exsolution
设单节电池电动势 $\mathcal E_0$。
每组串联后电动势
$\mathcal E_{\rm eq}=(N/n)\mathcal E_0$，
一组内阻 $(N/n)r$，
$n$ 组并联等效内阻
$R_{\rm int}=Nr/n^2$。
负载功率
\[
P_R=\frac{(N\mathcal E_0/n)^2 R}
{(R+Nr/n^2)^2}
=\frac{N^2\mathcal E_0^2Rn^2}
{(Rn^2+Nr)^2}.
\]
把 $n^2$ 视为连续变量，
求极大得 $Rn^2=Nr$，
所以
\[
\boxed{n_*=\sqrt{\frac{Nr}R}
=\sqrt{\frac{300(0.3)}{10}}=3}.
\]
\excheck $n=3$ 正好整除 $N=300$，
各并联组均有100节电池。

\electrocurrentproblem{199}{120}
图3.55中先将 $C=5.0\,\mathrm{\mu F}$
接电动势 $\mathcal E=200\,\mathrm V$
充电，随后开关由1拨到2，
经 $R_1=500\,\Omega$、
$R_2=330\,\Omega$ 放电。
求在 $R_1$ 中产生的焦耳热。
\exsolution
初态储能
$U_0=\frac12C\mathcal E^2$。
拨到2后电容通过
$R_1$、$R_2$ 的串联回路放电，
瞬时回路电流相同，
各电阻耗散功率比恒为
$P_1/P_2=R_1/R_2$。
因此总储能中由 $R_1$ 耗散的份额
\[
Q_1=U_0\frac{R_1}{R_1+R_2}
=\boxed{\frac{C\mathcal E^2R_1}{2(R_1+R_2)}}.
\]
代数值：
\[
\boxed{Q_1\simeq0.060\,\mathrm J=60\,\mathrm{mJ}}.
\]
\excheck 注意放电回路内两电阻\emph{串联}，
不是按照并联电阻分流比例分配热量。

\electrocurrentproblem{200}{120}
空气平板电容器原电容
$C=20\,\mathrm{nF}$，
接固定电压 $V=100\,\mathrm V$。
板间有金属插片，其厚度占间距
$\eta=0.60$，随后缓慢抽出。
求 (a) 电容器场能变化；(b) 外力机械功。
\exsolution
金属插片使有效电场间隙减为
$d(1-\eta)$，
插入时电容
$C_i=C/(1-\eta)$。
抽出后电容为 $C_f=C$，
电压始终由电源维持。
电容器储能 $U=\frac12CV^2$，
所以
\[
\boxed{\Delta U=
\frac12(C_f-C_i)V^2
=-\frac{CV^2\eta}{2(1-\eta)}}.
\]
恒电压下外力反抗吸引插片的电力，
外界机械功
\[
\boxed{A_{\rm ext}=-\frac12V^2(C_f-C_i)
=\frac{CV^2\eta}{2(1-\eta)}}.
\]
代入 SI 数值：
\[
\boxed{\Delta U=-0.15\,\mathrm{mJ},
\quad A_{\rm ext}=+0.15\,\mathrm{mJ}}.
\]
\excheck 电压由电源固定时，
电场储能变化与外界机械功不直接相等；
电源也交换电荷和能量。
